\documentclass[final]{beamer}

\usepackage[orientation=landscape,size=a0,scale=1.4]{beamerposter}
\usetheme{gemini}

\usepackage{amsmath,amssymb}
\usepackage[version=4]{mhchem}
\usepackage{booktabs}
\usepackage{siunitx}
\usepackage{graphicx}
\usepackage{tikz}
\usepackage{pgfplots}
\pgfplotsset{compat=1.18}

\newlength{\sepwidth}
\newlength{\colwidth}
\setlength{\sepwidth}{0.02\paperwidth}
\setlength{\colwidth}{0.30\paperwidth}
\newcommand{\separatorcolumn}{\begin{column}{\sepwidth}\end{column}}

\title{Quantum Transport in Topological Insulator Thin Films}
\author{Jane Doe \and John Smith \and Advisor Name}
\institute{Department of Physics, University of Example}

\begin{document}
\begin{frame}[t]
\begin{columns}[t]

\separatorcolumn
\begin{column}{\colwidth}

  \begin{block}{Introduction}
    Topological insulators (TIs) host metallic surface states protected by
    time-reversal symmetry, while the bulk remains insulating. These states
    are promising for spintronics and low-dissipation electronics.
    \begin{itemize}
      \item Bulk band gap: \SI{0.3}{\electronvolt}
      \item Dirac surface states with spin-momentum locking
      \item Robust against non-magnetic disorder
    \end{itemize}
  \end{block}

  \begin{block}{Theory}
    The low-energy surface Hamiltonian is
    \begin{equation}
      H = \hbar v_F (\sigma_x k_y - \sigma_y k_x),
    \end{equation}
    where $v_F \approx \SI{5e5}{\meter\per\second}$ and $\sigma_i$ are
    Pauli matrices. The resulting Dirac cone is shown in Fig.~\ref{fig:cone}.
  \end{block}

\end{column}

\separatorcolumn
\begin{column}{\colwidth}

  \begin{block}{Methods}
    Thin films (\SI{5}{}--\SI{20}{\nano\meter}) were grown by MBE on
    \ce{Al2O3}(0001). Transport measurements were performed in a
    dilution refrigerator down to \SI{20}{\milli\kelvin}.
  \end{block}

  \begin{block}{Results}
    \begin{figure}
      \centering
      \includegraphics{example-image-a}
      \caption{Longitudinal conductivity vs.\ temperature.}
      \label{fig:cone}
    \end{figure}
  \end{block}

\end{column}

\separatorcolumn
\begin{column}{\colwidth}

  \begin{block}{Discussion}
    The data are consistent with weak antilocalization from strong
    spin-orbit coupling. Deviations below \SI{50}{\milli\kelvin} suggest
    bulk leakage, which we attribute to inadvertent doping.
  \end{block}

  \begin{block}{Conclusions}
    \begin{enumerate}
      \item High-quality TI films grown by MBE.
      \item Surface-dominated transport above \SI{50}{\milli\kelvin}.
      \item Future work: gate-tuning the Fermi level.
    \end{enumerate}
  \end{block}

  \begin{block}{References}
    \small
    \begin{enumerate}
      \item M.~Z. Hasan and C.~L. Kane, Rev.\ Mod.\ Phys.\ \textbf{82}, 3045 (2010).
      \item X.-L. Qi and S.-C. Zhang, Rev.\ Mod.\ Phys.\ \textbf{83}, 1057 (2011).
    \end{enumerate}
  \end{block}

\end{column}
\separatorcolumn

\end{columns}
\end{frame}
\end{document}
